% =============================================================================
%  CHAPTER 27 -- DISCUSSION: WHAT THIS DOES AND DOES NOT REPLICATE
% =============================================================================
\section{Discussion: what this does and does not replicate}
\label{ch:nongoals}

A research prototype is defined as much by its non-goals as by its results. This chapter
states, for each thing a production statistical-arbitrage book has and this project does
not, what is missing, why it matters at fund level, which direction it would move the
results, and what it would take to add it. The ordering is by expected impact on the
conclusions.

\subsection{1. A survivorship-bias-free universe}

\textbf{Missing:} point-in-time index membership, delisted-ticker prices, and delisting
returns. \textbf{What we have:} 195 names that Yahoo returns a full history for today, with
a 2016-vintage membership table (Section~\ref{sec:data:survivorship}).

\textbf{Why it matters at fund level.} Survivorship bias is not a rounding error for
relative-value strategies; it removes precisely the tail events that kill them. A pair or
basket trade's worst outcome is a permanent relationship break --- an acquisition, a
bankruptcy, a business-model divergence --- and those outcomes are correlated with
delisting. Conditioning on names that survived to 2026 removes most of those events from
the sample. The measured MaxDD of $0.384$--$0.419$ and the CVaR figures are therefore
optimistic, and the Sharpe ratios are optimistic by an unknown but certainly non-zero
amount.

\textbf{Direction:} results would get \emph{worse}. \textbf{Fix:} join a yearly constituent
history (public sources exist) to a price source that retains delisted series (CRSP or a
vendor graveyard file), and apply delisting returns on the exit date. This is the single
highest-value data upgrade available and it is the only one that could plausibly change the
qualitative conclusions.

\subsection{2. Realistic financing and borrow costs}

\textbf{Missing:} stock-loan fees, borrow availability and recall risk, margin and
financing on the long leg, and the collateral yield on short proceeds.
\textbf{What we have:} a proportional 10\,bps one-way cost with a 0--50\,bps sweep, plus a
square-root impact model that is implemented but numerically irrelevant at \$1 notional.

\textbf{Why it matters.} A market-neutral book is a \emph{financed} position: the short
proceeds are collateralised, the long leg is levered, and the short leg pays a loan fee that
is 25--50\,bps/year on easy-to-borrow large caps and several percent on specials --- and
specials are exactly the names where mispricings are largest, so the cost is
adversarially correlated with the opportunity. Borrow recalls also force liquidation at the
worst moment, which is a path risk no proportional-cost model captures. At our measured
gross annual return of $3.02\%$ and $16.6$-bar average holding, a 30\,bps/year borrow on
half the notional costs $\approx0.5\%$/year --- enough to take the net return from $0.38\%$
to about zero.

\textbf{Direction:} worse. \textbf{Fix:} a per-name borrow-cost series (vendor data), a
recall-event model, and a financing term on the NAV recursion
\eqref{eq:bt:pnl}.

\subsection{3. Capacity}

\textbf{Missing:} any statement about how much money the strategy can hold.
\textbf{What we have:} a per-dollar edge on \$1 gross notional, with the impact model of
\eqref{eq:impact} implemented but non-binding at that scale.

\textbf{Why it matters.} Edge per dollar and dollars of edge are different products, and a
fund buys the second. Solving \eqref{eq:impact} for the participation rate at which impact
alone equals a 10\,bps budget ($\sigma=2\%$, $Y=0.5$) gives $Q/V\approx1\%$ of daily
volume --- of order \$15M/day traded per name in \code{ADBE} at its historical volume.
Across four baskets with $\approx12$ round trips per year, that implies a book capacity in
the low hundreds of millions before impact consumes the edge; the true limit is lower
because relative-value flow is concentrated in the same names at the same times.

\textbf{Direction:} unchanged for the reported ratios, but it bounds the economic value of
any improvement. \textbf{Fix:} re-run the backtest with notional scaled to a participation
constraint, and report net Sharpe as a function of assets under management --- the standard
capacity curve.

\subsection{4. A real risk model and factor neutrality}

\textbf{Missing:} neutralisation against a multi-factor risk model (market, size, value,
momentum, sector, volatility, liquidity), position limits by sector and by name,
ex-ante tracking-error budgets, and a leverage policy.
\textbf{What we have:} dollar neutrality by construction
\eqref{eq:grossing}, inverse-vol strategy weights with a 45\% cap and verified gross/net
limits (Section~\ref{sec:bt:sizing}), and realized VaR/CVaR.

\textbf{Why it matters.} A book that is dollar-neutral is not risk-neutral. The
Phase~\ref{ch:phase7} regime decomposition shows the losses concentrate in the 2022
rate-driven rotation --- a factor event, not an idiosyncratic one --- which means the book
carried an unneutralised growth/duration exposure that the dollar-neutral construction did
not see. A fund would size this book against a factor model and would have reduced or
hedged that exposure, changing both the return and the drawdown.

\textbf{Direction:} ambiguous --- probably lower return \emph{and} lower drawdown, i.e.\ a
better Sharpe but a smaller absolute P\&L. \textbf{Fix:} regress each strategy's daily P\&L
on a public factor set, neutralise the residual betas in the position construction, and
report ex-ante factor exposures alongside the realized tail statistics.

\subsection{5. Intraday and higher-frequency implementation}

\textbf{Missing:} anything below daily bars. No intraday data, no order-level execution, no
queue position, no partial fills.
\textbf{What we have:} daily closes with next-bar execution, which is the most conservative
realistic timing at daily frequency.

\textbf{Why it matters.} Chapter~\ref{ch:multipletesting} shows the binding constraint on
this entire research programme is \emph{statistical power}: $T^\ast=(z/\mathrm{SR}_d)^2$
days. Moving to hourly bars multiplies $T$ per year by $\approx7$ and, more importantly,
shortens half-lives from $\approx9.4$ days to hours, which raises the number of independent
bets per year by one to two orders of magnitude. By the fundamental law
\eqref{eq:fundamentallaw} that is where a weak per-bet IC becomes a testable book Sharpe.

\textbf{Direction:} this is the one change that could plausibly turn the negative result
positive --- and also the one that changes the problem entirely (microstructure costs,
latency competition, queueing). \textbf{Fix:} an intraday data source and an execution
simulator with fill modelling; the C++ engine's per-bar cost ($\approx600$\,ns) already
supports $10^6$-bar simulations comfortably.

\subsection{6. Breadth: hundreds of baskets instead of four}

\textbf{Missing:} a systematic basket-construction rule applied across the whole universe.
\textbf{What we have:} four hand-specified economically coherent baskets, chosen a priori by
sub-industry.

\textbf{Why it matters.} With four baskets the book has $\approx12$--$50$ independent bets a
year and, per \eqref{eq:sharpetstat}, no chance of significance. With 200 disjoint baskets
the book-level Sharpe of the same per-basket edge would be $\approx\sqrt{50}\times$ larger in
$t$-statistic terms, and the relevant test becomes one test on the book rather than 200 tests
on its parts. This is the cheapest path to a falsifiable experiment and it is why the C++
engine's speed was worth building for.

\textbf{Direction:} better power, and possibly a real (small) positive book Sharpe --- or a
definitive negative, which is equally valuable. \textbf{Fix:} an automated
basket-construction rule (sub-industry peers by market cap, or lasso-selected from the
sector), run over the 195-name universe, with the full search audited under
Chapter~\ref{ch:multipletesting}'s machinery and Benjamini--Hochberg FDR control rather than
Bonferroni (Section~\ref{sec:mt:bonferroni}).

\subsection{7. Alternative data and richer features}

\textbf{Missing:} fundamentals, options-implied quantities, short-interest, ownership,
news/sentiment, and any non-price input. \textbf{What we have:} 11 causal price-and-volume
features (Section~\ref{sec:nonlinear:target}).

\textbf{Why it matters.} The nonlinear verdict of Phase~\ref{ch:phase9} --- ``linear wins''
--- is a statement about a feature set with little nonlinear structure to find. Short
interest and borrow cost, options skew, and earnings-date proximity all plausibly interact
nonlinearly with reversion speed, and they are the features that would make a GBM worth its
variance. Without them, the honest conclusion is that nonlinearity is not detectable
\emph{in prices alone at daily frequency}, not that it is absent.

\textbf{Direction:} unknown; potentially much better, and the only route to a genuinely
higher IC. \textbf{Fix:} a separate data infrastructure (vendors, licensing, point-in-time
storage to avoid look-ahead in fundamentals restatements) --- a project of comparable size to
this one.

\subsection{8. Cross-asset integration}

\textbf{Missing:} anything outside U.S.\ large-cap equities --- no rates, FX, credit,
commodities, or futures; no cross-asset relative value; no hedging of the book's residual
rate exposure with, e.g., Treasury futures.

\textbf{Why it matters.} The 2022 loss (net Sharpe $-1.306$ in the selloff window) was
driven by a macro factor. A cross-asset overlay would either hedge it or, better, monetise
it: equity-bond and within-complex relative value in a rate shock is a different and
complementary strategy. Real books are assembled across asset classes precisely because
their stress periods differ.

\textbf{Direction:} better tail behaviour, more diversification, more complexity.
\textbf{Fix:} a distinct research programme; the engine's data layer would need futures roll
and multi-currency handling.

\subsection{9. Live capital allocation and operational feedback}

\textbf{Missing:} everything about running money --- a live risk limit framework with
kill-switches, prime-brokerage and clearing integration, position and cash reconciliation,
corporate-action handling in real time, compliance and best-execution obligations, model
monitoring and drift detection, and the human process that decides when to stop a strategy.

\textbf{Why it matters.} A backtest's P\&L is a deterministic function of history; a live
book's P\&L is the output of a control system with feedback, delays, and operational
failure modes. The most common way a statistically sound relative-value strategy loses money
in production is not that the signal decays but that a break goes undetected, a limit is
breached during a data outage, or a position cannot be exited when the borrow is recalled.
None of that is in scope here and none of it is captured by any number in this book.

\textbf{Direction:} worse, always. \textbf{Fix:} an operational build, not a research one.
The one piece that \emph{is} research and should be added first is drift detection: a
monitoring statistic (e.g.\ a rolling DM test of live P\&L against backtest expectation) that
triggers a review when the realized behaviour diverges from the model's.

\subsection{10. Colocated, FPGA-level execution}

\textbf{Missing:} any latency competition. \textbf{What we have:} a
$\approx600$\,ns-per-bar simulation core on a shared 2-core sandbox, with per-stage
percentiles measured to $\pm15.8\%$ bucket resolution.

\textbf{Why it matters --- and why it does not, here.} Colocation and FPGA matter for
strategies whose edge decays in microseconds. Ours decays over $\approx9.4$ \emph{days}: the
half-life is eight orders of magnitude longer than the latency budget that colocation buys.
The engine's low-latency design is therefore not an alpha claim; it is (i) a specification
deliverable, (ii) the enabler of the breadth-first research programme of item 6 (a
$10^6$-backtest scan), and (iii) an honest statement that the simulation cost is not the
constraint on any conclusion in this book. Claiming otherwise would be the sort of
over-claim this chapter exists to prevent.

\subsection{Summary table of non-goals}

\begin{center}
\renewcommand{\arraystretch}{1.25}
\small
\begin{tabular}{p{0.30\linewidth}p{0.20\linewidth}p{0.40\linewidth}}
\toprule
Non-goal & Direction of bias & Cheapest credible fix \\
\midrule
Survivorship-bias-free universe & optimistic & historical membership + delisted prices \\
Borrow / financing costs & optimistic & per-name loan-fee series + financing term in \eqref{eq:bt:pnl} \\
Capacity analysis & silent & participation-constrained re-run, Sharpe vs AUM curve \\
Multi-factor neutralisation & ambiguous & public factor set, neutralise in position construction \\
Intraday frequency & \emph{could help} & intraday data + fill model; engine already fast enough \\
Universe-wide breadth & \emph{could help} & automated basket rule over 195 names, FDR-controlled audit \\
Alternative data & \emph{could help} & separate data infrastructure (point-in-time) \\
Cross-asset integration & better tails & separate research programme \\
Live operational loop & worse & not a research task; add drift monitoring first \\
Colocation / FPGA & irrelevant at this half-life & none needed --- and none claimed \\
\bottomrule
\end{tabular}
\end{center}
